\begin{table}[htbp]
\centering
\caption{Result 1: The Border Premium in Revenue per Dispensary}
\label{tab:result1}
\small
\begin{tabular}{p{6.6cm}cccc}
\toprule
Specification & Coef. & SE (HC1) & SE (cl.) & $p$ (cl.) \\
\midrule
\multicolumn{5}{l}{\textit{Binary: border ($\le$70mi) vs.\ interior}} \\
No controls & $+0.382$ & $0.099$ & $0.219$ & $0.080$ \\
Full controls$^{a}$ & $+0.341$ & $0.075$ & $0.161$ & $0.033$ \\
\quad excl.\ Sedgwick County, CO & $+0.302$ & $0.068$ & $0.147$ & $0.040$ \\
Colorado only & $+0.552$ & -- & $0.199$ & $0.005$ \\
Washington only & $+0.233$ & -- & $0.243$ & $0.334$ \\
\midrule
\multicolumn{5}{l}{\textit{Continuous: log(distance to nonlegal border, miles)}} \\
Full controls$^{a}$ & $-0.215$ & $0.048$ & $0.095$ & $0.023$ \\
\quad excl.\ Sedgwick County, CO & $-0.168$ & $0.048$ & $0.092$ & $0.067$ \\
\bottomrule
\end{tabular}
\begin{minipage}{0.95\textwidth}
\vspace{0.3em}
\footnotesize
\textit{Notes:} Outcome is log(revenue per active dispensary). County-year panel, Colorado and Washington,
$n=787$ (80 counties: 42 Colorado, 38 Washington; 12 border counties: 6 Colorado, 6 Washington; 68 interior).
$^{a}$Full controls: log(population), year fixed effects, state fixed effects. Both the binary and
continuous measures are significant under county-clustered inference; a county closer to a nonlegal border
earns significantly more revenue per dispensary, and the relationship holds using a continuous distance
measure rather than only the 70-mile cutoff.
\end{minipage}
\end{table}
